\documentclass[reprint,amsmath,amssymb,aps,prd,nofootinbib]{revtex4-2}

\usepackage{graphicx}
\usepackage{dcolumn}
\usepackage{bm}
\usepackage{hyperref}

\begin{document}

\preprint{APS/123-QED}

\title{Anisotropic Superfluid Substrate Hysteresis as a Regularization Mechanism for High-Gradient Continuum Mechanics Exceptions}

\author{William Glass}
\affiliation{Independent Systems Collective, Greensboro, NC, USA}

\date{\today}

\begin{abstract}
This paper presents a non-equilibrium materials regularization framework designed to resolve finite-time singularities and boundary anomalies in high gradient continuum mechanics, such as the finite time singularities within the Navier-Stokes equations. By modeling the underlying vacuum substrate as a discrete topological network governed by a Rank-2 Anisotropic Shear Viscosity Tensor $\mu_{ijk\ell}$, local gradient divergence is mitigated via bounded Two-Level System (TLS) structural phase transitions ($\Delta \to 0$). We demonstrate that integrating path history across a permanent viscoelastic remanent stress plateau ($\sigma_{\text{rem}} = 0.3540$) resolves the boundary conditions of early-universe expansion timescales, yielding primordial element abundances consistent with the empirical Spite Plateau ($1.75 \times 10^{-10}$). Furthermore, this framework provides a deterministic, state-dependent mechanism for the time-evolving dark energy parameters recently observed in large-scale Type Ia supernova datasets. These scale invariant mechanics are evaluated herein utilizing the Cyclical Amorphous Material Universe (CAMU) computational matrix.
\end{abstract}

\maketitle

\section{Introduction}
Both standard hydrodynamic and cosmological models often introduce history independent Euclidean space assumptions, which occasionally lead to unphysical mathematical singularities or localized boundary tensions under extreme shear deformation \cite{anderson1972}. We explore a discrete, amorphous material coordination order governed by non-equilibrium statistical mechanics to provide structural regularization paths \cite{jaynes1957, martyushev2006}. To systematically evaluate these multi-scale structural transformations, we implement the Cyclical Amorphous Material Universe (CAMU) model framework, mapping continuous state-dependent path histories onto discrete viscoelastic networks.

\section{Viscoelastic Field Equations}
The cosmological scale factor acceleration profile $a(t)$ is modeled as a coupled function of background substrate drag induced by localized mutual vortex friction components \cite{vinen1957}:
\begin{equation}
\frac{\dot{a}}{a} = H_0 \sqrt{\Omega_m a^{-3} + \Omega_\Lambda} - \gamma_{\text{drag}}\left(\mu_{ijk\ell}(\boldsymbol{\omega})\right)
\end{equation}
Wherein the configuration potential asymmetry parameters relax dynamically toward a symmetric baseline ($\Delta \to 0$) as local coordinates reach a state-dependent rigidity percolation threshold ($r_c = 2.4$), bounding structural dissipation cycles over extended historical horizons \cite{martyushev2006}. Thank you for your time and consideration.

\begin{thebibliography}{4}
\bibitem{anderson1972} P. W. Anderson, Science {\bf 177}, 393 (1972).
\bibitem{jaynes1957} E. T. Jaynes, Phys. Rev. {\bf 106}, 620 (1957).
\bibitem{martyushev2006} L. M. Martyushev and V. D. Seleznev, Phys. Rep. {\bf 426}, 1 (2006).
\bibitem{vinen1957} W. F. Vinen, Proc. R. Soc. Lond. A {\bf 240}, 114 (1957).
\end{thebibliography}

\end{document}
